% Immutable cited-input supplement. Target proofs remain embargoed.
% Target: push:src:E.2.7; tag 033F.
% Frozen source SHA256 74497bfe327fae25c9cdcb17334b3bf42a9793a9df7c0a4404d97568f2f760f0.
% CompendiumRef is resolved against the cited S6_09 source edition, not silently identified with the current series statement.

% Exact cited input: series:I / push:src:E.2.1
\begin{theorem}[Tangent cohomology, direct images and miniversality]\hypertarget{tag-032U}{}\label{push:src:E.2.1}
\hypertarget{tag-01RH}{}\label{fam:src:A.5.2}\label{fam:st:A.5.2}\label{inv:D1.3}\label{inv:D2.1}\label{inv:D3.1}\label{cdp:st:E9}
\hypertarget{tag-02YO}{}\label{night:P5.3}\label{stmt:P5.3}
\hypertarget{tag-02YP}{}\label{night:P5.4}\label{stmt:P5.4}

Let $k,l$ be positive integers with $\gcd(3k,4l)=1$, and let $X$
be a sphere member with multiplicities $(3k,4l)$ and the admissible
logarithmic data of \cref{push:section:E.1}. Put $z=kl$ and
$d=3k+4l$. On $\mathfrak B_{k,l}=\mathbb P(3k,4l)$, the vertical
root bundle $E$, the root lift $\bar f$ and its direct images satisfy
\[
 \begin{split}
 E&=\mathcal O\oplus\mathcal O(-z),\\
 R^q\bar f_*\mathcal O_X&=
 \mathcal O,\quad\mathcal O\oplus\mathcal O(-z),\quad\mathcal O(-z),\\
 K_X&=\bar f^*\mathcal O(z-d).
 \end{split}
\]
The tangent cohomology is
$\bigl(h^0,h^1,h^2,h^3\bigr)(X,T_X)=(1,1,1,0)$.
For the transition-arrow convention
$\kappa(\Phi)=(\partial_c\Phi)\Phi^{-1}$, the Kodaira--Spencer
class of $\partial_c$ is $e\otimes[\gamma]$, and the class of
$\partial_u$ is $-(2\pi i u)^{-1}e\otimes[\gamma]$.
Both classes are nonzero. The primary Kuranishi bracket vanishes,
the Kuranishi germ is a smooth reduced disc, and the fixed-data
period family is locally miniversal.

For the standard $(3,4)$ data, $e=\xi$ and the coarse direct
images are
\begin{equation}
 f_*T_X=\mathcal O\oplus\mathcal O(-1),\qquad
 R^1f_*T_X=\mathcal O\oplus\mathcal O(-1)^2,\qquad
 R^2f_*T_X=\mathcal O.
\label{pub:I:eq:D-4}
\end{equation}
\end{theorem}

% Exact cited input: series:I / push:src:E.1.2
\begin{theorem}[Classification in Viaclovsky's coordinates]\hypertarget{tag-032S}{}\label{push:src:E.1.2}


Use a local logarithm $c$ of the transported linearization,
$u=e^{-2\pi i c}$, and the same reference section and cusp comparison
on both sides.  Two sphere members with data
$(m,\rho_0,N_0;n,\rho_1,N_1;c)$ and primed data are biholomorphic
if and only if their multiplicities agree and there are
$\alpha\in\{1,-1\}$ and $h=c'-c\in\mathbb Z$ such that
\[
 b'=\alpha b,\qquad
 hN'_0\equiv\rho'_0-\alpha\rho_0\pmod{k},\qquad
 hN'_1\equiv\rho'_1-\alpha\rho_1\pmod{2l}.
\]
Every such biholomorphism covers the identity of $B$.
For fixed discrete data the period identification subgroup is
$2kl\mathbb Z$.  At fixed $u$ all sphere translation data give
exactly $2kl$ classes.  A positive $u$-circuit acts on them by a
single cycle.
\end{theorem}

% Exact cited input: series:I / push:src:E.1.3
\begin{corollary}[Holomorphic dependence of the comparisons]\hypertarget{tag-032T}{}\label{push:src:E.1.3}


For fixed discrete data in \ref{push:src:E.1.2} and a fixed admissible shift $h$,
the constructed comparisons are holomorphic isomorphisms of families
over a sufficiently small parameter disc.  In particular, for sphere
data with multiplicities $(3,4)$ the identification with the standard
family uses the parameter
\[
 W=e^{-\pi i(c+s_2)}=(-1)^{s_2}w.
\]
The smooth monodromy conjugacy class around the positive parameter
circle is unchanged by this identification of twist data.
\end{corollary}

% Exact cited input: series:I / push:src:B.4.4
\begin{theorem}[Extension of the marked base map]\hypertarget{tag-02RE}{}\label{night:P2.4}\label{stmt:P2.4}

\hypertarget{tag-0314}{}\label{push:src:B.4.4}



For every admissible $(3,4)$ twist, the marked map $f:X_c\to\Pp^1$ extends meromorphically over every proper Hausdorff analytic filling of a finite parameter cover. For a smooth given total space, its graph is a proper projective modification carrying the marked map. In general, a smooth proper modification carries a holomorphic marked map to $\Pp^1$, and hence a proper map to $\Pp^1\times\Delta$.
\end{theorem}

